\begin{comment}

Adicionar:
  - adicionar imagem do teorama CAP
  - adcionar imagem do resultado de impossibilidade FLP

Melhorar:
  - Fluidez entre paragráfos

Estrutura do conteúdo:
  - Explicar a relação entre comunicação assíncrona e sistemas Offline-First
  - Ressaltar o cenário ideal onde seria possível manter consistência e disponibilidade com tolerância a falhas
  - Mostrar como um sistema completamente assíncrono atende a essa demanda
  - Discorrer sobre o Teorema CAP para entender como não é possível manter consistência e disponibilidade com tolerância a falhas
  - Discorrer sobre como não é possível ter um sistema completamente assíncrono de acordo com o FLP
  - Concluir como essas ideias justificam a existência de estratégias de sincronização de dados

\end{comment}

\subsubsection{TEOREMA \en{CAP}}

Teorema \en{CAP} define um \en{tradeoff} entre consistência e disponibilidade em sistemas distribuídos não confiáveis. Apresentado por Eric Brewer através de um exemplo generalista de serviço \en{web} para demonstrar essa troca entre disponibilidade, consistência e tolerância a partição. Consistência sendo a capacidade do serviço de retornar a resposta correta para uma requisição; Disponibilidade é a garantia de toda requisição receber uma resposta, mesmo que eventualmente; e, tolerância a partição, este é a capacidade serviço de permanecer funcionando mesmo no caso de perda de comunicação \cite{CAPtheorem2012}.

No contexto de restrições de sistemas distribuídos, pode se destacar, ainda, o resultado de impossibilidade \en{FLP}; este afirma que em um sistema completamente assíncrono, isso é, em um sistema onde não há garantia de quando uma menssagem será entregue, não existe algoritmo determinístico capaz de encontrar consenso desde que o sistema assuma pelo menos uma falha de execução. Desse modo, entendendo o contexto com conectividade intermitente, é necessário, na construção de sistemas reais, relaxar a assincronicidade \cite{FLP1985}, por exemplo, ao aplicar mecanismos como: consistência eventual forte \cite{Gomes2017VerifyingSEC} ou \en{CRDTs}.

Assim, na construção da aplicação casos onde ocorrem desconexão de par devem ser resolvidos com aplicação de estratégias para manter os dados estruturados e coesos.
